\documentclass[instructions.tex]{subfiles}
\begin{document}
 
\chapter{Plusieurs croient pardonner, sans pardonner en effet.}
 
\primo Je lui pardonne pour cette fois, mais qu’il ne s’y retrouve plus. Ce n’est pas une fois, mais septante fois sept fois, dit Jésus-Christ, que vous devez pardonner, c’est-à-dire toujours. Vous avez toujours besoin de la miséricorde de Dieu; vous devez donc toujours exercer votre miséricorde, et l’exercer envers tous.

\secundo Je lui pardonne, mais j’ai bonne mémoire; jamais je n’oublierai ce qu’il m’a fait. C’est donc à dire que vous ne voulez pas que Dieu oublie vos péchés. Si vous ne pouvez oublier l’injure et le tort qu’on vous a faits, il vous est permis d’y penser. Mais comment? c’est en adorant les desseins de Dieu, qui vous éprouve; c’est en regardant votre ennemi comme l’instrument dont Dieu se sert pour travailler à votre couronne. Semëis, vomissant des injures, et lançant des pierres contre David, son roi : Laissez-le faire, disait ce saint roi, Dieu le permet pour m’éprouver; Imitez cet exemple, ne pensez aux injures et aux torts qu’on vous fait, que pour bénir le Seigneur de tout ce qui vous arrive \footnote{Benedicam Dominum in omni tempore. P.33}.

\tertio Je lui pardonne, mais je ne veux pas le voir. Ah! malheureux, Dieu ne vous voit-il pas, quelque indigne que vous soyez? comment osez-vous prier le Seigneur de vous regarder d’un œil de miséricorde, si vous ne voulez pas voir votre frère? Si vous n’aimez comme vous devriez un œil, disait saint François de Sales, je le regarderais toujours avec l’autre de bon œil et de bon cœur.

Si une personne était scandaleuse, et que sa compagnie fût pernicieuse à votre âme, dans ce cas il vous est permis de l’éviter avec prudence et précaution, mais non pas par rancune ni par orgueil: fuyez sa compagnie, si elle vous est funeste, mais aimez sa personne en vue de Dieu.

Lorsque vous avez de l’autorité sur une personne, il vous est quelquefois permis de lui témoigner de l’indifférence ou de la froideur, pour lui faire sentir sa faute et son devoir; mais que ce soit par le motif d’une charitable correction, et qu’il n’y ait ni haine, ni aversion de votre part, ni scandale du côté du public.

\quarto Mais mon ennemi ne veut ni me voir, ni me parler. Qu’en savez-vous? Peut-être le désire-t-il plus que vous. S’il ne veut pas vous recevoir, n’est-ce point parce que, loin de lui parler avec honnêteté, vous ne lui parlez qu’avec hauteur, en lui faisant des reproches? Ce n’est pas ainsi qu’on se réconcilie; on n’éteint pas l’embrasement en jettant l’huile dans le feu. Si vous craignez de l’irriter et de vous compromettre, en paraissant devant lui, priez un ami commun de lui parler de votre part, pour lui faire agréer votre visite. Cette démarche vous fera honneur.

Je lui pardonne, mais je le laisse tel qu’il est. Vous voulez donc que Dieu vous laisse tel que vous êtes; c’est-à-dire, que vous voulez que Dieu vous réprouve; car si Dieu vous laissait à vous-même, vous seriez un réprouvé. Pardonner de la sorte, ce n’est point pardonner.

Je lui veux point de mal. Ce n’est pas assez; il faut lui vouloir du bien, l’aimer comme vous-même, être affligé, si on lui fait du mal, l’empêcher dans l’occasion, etc. Vous dites que vous ne lui voulez pas de mal; cependant, s’il se présente une occasion de lui rendre un mauvais office, vous la saisissez; si on le calomnie, vous y prenez plaisir. Vous portez pendant dix ans, dans le cœur, l’injure qu’il vous a faite; vous en parlez, vous vous en expliquez durement avec vos amis et devant vos enfans. C’est ainsi qu’on transmet, qu’on éternise les rancunes dans les familles, et qu’en disant froidement : Je ne veux point de mal, on a toujours le fiel dans l’âme. On sait ainsi, au scandale du public, allier la communion avec le ressentiment. Communion abominable, dit Tertullien, qui, n’ayant pas la force de réunir les hommes ensemble, ne peut les réconcilier avec Dieu.
 
\section{Résolutions}
 
\primo Je pardonnerai toujours, je verrai de bon œil mon ennemi; je lui ferai du bien dans l’occasion, et je prendrai tous les moyens que la charité et la prudence m’indiqueront pour me réconcilier promptement et d’une manière sincere avec lui. 

\prayer{Mon Dieu, pour accomplir ces résolutions, j’ai besoin du secours de votre grâce, d’une profonde humilité, d’un grand détachement de moi-même, et de beaucoup de vigilance sur mon propre cœur; venez donc puissamment à mon secours.}
 
\end{document}
